% Part: first-order-logic
% Chapter: axiomatic-deduction
% Section: proof-theoretic-notions

\documentclass[../../../include/open-logic-section]{subfiles}

\begin{document}

\iftag{FOL}
      {\olfileid[zh]{fol}{axd}{ptn}}
      {\olfileid[zh]{pl}{axd}{ptn}}

\olsection{证明论诸概念}

\begin{explain}
正如我们定义了若干重要的语义概念
(\iftag{FOL}{有效性}{重言式}、衍推、可满足性)，我们现在
定义相应的\emph{证明论概念}。这些概念并不是
通过诉诸!!{sentence}在!!{structure}中的满足来定义的，
而是通过诉诸某些公式的!!{derivability}或!!{nonderivability}来定义的。
这些概念相一致乃是一项重要发现。
它们相一致正是\emph{可靠性}定理与\emph{完全性}定理的内容。
\end{explain}

\begin{defn}[!!^{derivability}]
若存在从~$\Gamma$ 出发并终于~$!A$ 的!!a{derivation}，则!!^a{formula}~$!A$ 由 $\Gamma$ 是 \emph{!!{derivable}}，
记为 $\Gamma \Proves !A$。
\end{defn}

\begin{defn}[定理]
若存在从空集出发的对 $!A$ 的!!a{derivation}，则!!^a{formula}~$!A$ 是\emph{定理}。若 $!A$ 是定理，我们记
$\Proves !A$；若不是则记 $\Proves/ !A$。
\end{defn}

\begin{defn}[一致性]
由!!{formula}组成的集合 $\Gamma$ 是\emph{一致的}，当且仅当
$\Gamma\Proves/ \lfalse$；否则它是\emph{不一致的}。
\end{defn}

\begin{prop}[自反性]
\ollabel{prop:reflexivity}
若 $!A \in \Gamma$，则 $\Gamma \Proves !A$。
\end{prop}

\begin{proof}
  !!{formula}~$!A$ 本身即是从~$\Gamma$ 出发对~$!A$ 的!!a{derivation}。
\end{proof}

\begin{prop}[单调性]
\ollabel{prop:monotonicity}
若 $\Gamma \subseteq \Delta$ 且 $\Gamma \Proves !A$，则 $\Delta
\Proves !A$。
\end{prop}

\begin{proof}
从 $\Gamma$ 出发对 $!A$ 的任何!!{derivation}也是从~$\Delta$ 出发对
$!A$ 的!!a{derivation}。
\end{proof}

\begin{prop}[传递性]
\ollabel{prop:transitivity}
若 $\Gamma \Proves !A$ 且 $\{!A\} \cup \Delta \Proves
!B$，则 $\Gamma \cup \Delta \Proves !B$。
\end{prop}

\begin{proof}
  设 $\{!A\} \cup \Delta \Proves !B$。则存在从~$\{!A\} \cup
  \Delta$ 出发的!!a{derivation} $!B_1$, \dots, $!B_l = !B$。该!!{derivation}中的
  某些步骤之所以正确，是因为所用的规则引用了先前的一行~$!B_i = !A$。由
  假设，存在从~$\Gamma$ 出发对~$!A$ 的!!a{derivation}，即!!a{derivation}~$!A_1$, \dots, $!A_k = !A$，其中每个 $!A_i$ 或是一条公理，
  或是~$\Gamma$ 的!!a{element}，或由某条推理规则正确得出。现在考虑如下序列
  \[
  !A_1, \dots, !A_k = !A, !B_1, \dots, !B_l = !B.
  \]
  这是从 $\Gamma \cup \Delta$ 出发对~$!B$ 的一个正确的!!{derivation}，
  因为每个 $B_i = !A$ 现在都由同一规则证成，即证成~$!A_k = !A$ 的那条规则。
\end{proof}

注意，这表明特别地，若 $\Gamma \Proves !A$ 且 $!A
\Proves !B$，则 $\Gamma \Proves !B$。由此亦可得，若 $!A_1,
\dots, !A_n \Proves !B$ 且对每个~$i$ 有 $\Gamma \Proves !A_i$，则
$\Gamma \Proves !B$。

\begin{prop}
\ollabel{prop:incons}
$\Gamma$ 是不一致的当且仅当对每个~$!A$ 有 $\Gamma \Proves !A$。
\end{prop}

\begin{proof}
习题。
\end{proof}

\tagprob{FOL}
\begin{prob}
证明 \olref[fol][axd][ptn]{prop:incons}。
\end{prob}
\tagendprob

\tagprob{notFOL}
\begin{prob}
证明 \olref[pl][axd][ptn]{prop:incons}。
\end{prob}
\tagendprob

\begin{prop}[紧致性]
\ollabel{prop:proves-compact}
  \begin{enumerate}
  \item 若 $\Gamma \Proves !A$，则存在有穷子集 $\Gamma_0
    \subseteq \Gamma$ 使得 $\Gamma_0 \Proves !A$。
  \item 若~$\Gamma$ 的每个有穷子集都
    一致，则 $\Gamma$ 一致。
  \end{enumerate}
\end{prop}

\begin{proof}
  \begin{enumerate}
    \item 若 $\Gamma \Proves !A$，则存在!!{formula}的一个有穷序列
      $!A_1$, \dots,~$!A_n$ 使得 $!A \ident !A_n$，并且
      每个 $!A_i$ 或是一条逻辑公理，或是~$\Gamma$ 的!!a{element}，
      或由分离规则从先前的!!{formula}推出。取
      $\Gamma_0$ 为那些属于~$\Gamma$ 的 $!A_i$。则此!!{derivation}同样是从~$\Gamma_0$ 出发的!!a{derivation}，
      故 $\Gamma_0 \Proves !A$。
    \item 这是~(1) 在特殊情形 $!A
      \ident \lfalse$ 下的逆否命题。
\end{enumerate}
\end{proof}

\end{document}
